\begin{abstract}
Spiking Transformers expose binary internal states, but it remains unclear
whether their query--key patterns form reliable and compact lookup-table (LUT)
addresses. We study this question in QKFormer through a negative-controlled
pipeline covering address feasibility, held-out response reconstruction,
support-aware fallback, and subspace decomposition. Correct Q/K-addressed
prototypes reduce held-out response MSE by $3.71\%$ on CIFAR-10 and $4.42\%$
on CIFAR-100, compared with $2.79\%$ and $3.42\%$ for coarse token/channel
lookup. Shuffling the address--prototype alignment instead degrades MSE by
$17.83\%$ and $15.82\%$, showing that table capacity alone does not explain the
effect. Existing checkpoint and calibration-subset studies preserve this
ordering. To avoid a monolithic address table, we introduce a
subspace-decoupled estimator that composes token/channel, Q/gate, and K
residual subtables. In the canonical CIFAR-100 seed-42 analysis, TC+Q/gate and
TC+QK retain $87.58\%$ and $94.31\%$ of full-address reconstruction gain while
using $8.83\%$ and $20.54\%$ of the compact address-space entry count. Their
matched shuffled-subspace controls are worse than the global mean. These
results support compact, interpretable Q/K lookup-addressability for local
response reconstruction. In an end-to-end current-replacement diagnostic,
calibrated Q/K lookup preserves QKFormer accuracy within $0.44$ percentage
points on CIFAR-10/100 at $T=1$ and $T=4$, including a two-stage CIFAR-10
replacement within $0.27$ points. For all four attention projections, an exact
2-bit grouped spike-pattern LUT preserves CIFAR-100 accuracy within $0.06$
points at both $T=1$ and $T=4$, using 2,664 KiB of FP32 table values.
In a stricter native-operator audit, all 32 Conv/Linear affine operators,
their 31 paired inference BN modules, and all 35 LIF state transitions are
removed from the evaluation graph and replaced by LUT-native operators,
preserving CIFAR-100 accuracy within $0.20$ and $0.22$ points at $T=1$ and
$T=4$, respectively.
Across five fixed calibration subsets, aligned current lookup ranks first in
all five runs; its paired advantages over global and token/channel controls
have 95\% intervals that exclude zero. To test transfer beyond the native
QKFormer backbone, we evaluate Spikformer variants with and without the
QK-compatible interaction contract. The unmodified SSA geometry is retained as
a compact mismatch probe, where temporal+channel calibration recovers static
replacement from $66.47\%$ to $76.35\%$. More importantly, adapting the SSA
interaction to a QK-compatible contract yields $86.16\%$ calibrated replacement
from an $87.54\%$ clean baseline and exceeds a matched shuffled control by
$0.33$ points. These results support bounded, contract-conditioned transfer for
QK-compatible spiking-attention contracts. Broader Spiking Transformer and
event-data generalization is being evaluated separately, and the present
evidence does not support measured SRAM, latency, energy, ImageNet, or
production full-graph hardware gains.
\end{abstract}
